%! Parametrically Defined Circles and Lines — Unit 8, Lesson 8.3 — Warm-Up (Answer Key)
\documentclass[10pt]{article}
\usepackage{precalculus-article}
\usepackage{precalculus-key}   % pulls in -boxes; do NOT also load -boxes
\newcommand{\TallMath}[1]{$\displaystyle #1\rule[-1.4em]{0pt}{3.2em}$}

\begin{document}

\pageheader{Unit 8, Lesson 8.3}{Warm-Up --- Answer Key}
\namedateperiod

\vspace{0.10in}

\begin{notesbox}{Spiral Review --- Key}

\textbf{1.}\quad Complete the table of exact values.

\vspace{4pt}
\begin{center}
\renewcommand{\arraystretch}{1.6}
\begin{tabular}{|c|c|c|c|c|c|}
\hline
$t$ & $0$ & $\dfrac{\pi}{2}$ & $\pi$ & $\dfrac{3\pi}{2}$ & $2\pi$ \\[6pt]
\hline
$\cos t$ & \ans{$1$} & \ans{$0$} & \ans{$-1$} & \ans{$0$} & \ans{$1$} \\[4pt]
\hline
$\sin t$ & \ans{$0$} & \ans{$1$} & \ans{$0$} & \ans{$-1$} & \ans{$0$} \\[4pt]
\hline
\end{tabular}
\end{center}

\vspace{0.14in}
\textbf{2.}\quad The graph of $y = \sin(x)$ is shifted \textbf{up 3} units and
\textbf{scaled vertically by 2}. Write the equation of the transformed function.

\vspace{4pt}
\ans{$y = 2\sin(x) + 3$}

\vspace{0.12in}
\textbf{3.}\quad A point starts at $(1,\,0)$ and travels \emph{counterclockwise}
around a circle of radius 1. Fill in the table.

\vspace{4pt}
\begin{center}
\renewcommand{\arraystretch}{1.6}
\begin{tabular}{|l|c|}
\hline
\textbf{Turn} & \textbf{New coordinates} \\
\hline
After a quarter turn ($\tfrac{1}{4}$ of the way around) & \ans{$(0,\,1)$} \\
\hline
After a half turn ($\tfrac{1}{2}$ of the way around)    & \ans{$(-1,\,0)$} \\
\hline
After a full turn (all the way around)                  & \ans{$(1,\,0)$} \\
\hline
\end{tabular}
\end{center}

\begin{teachernote}
  Q1 activates the cosine and sine values students need for the unit-circle
  parametrization. Q2 previews the transformation template $(h + r\cos t,\;
  k + r\sin t)$ --- if students struggle here, connect the vertical shift to
  the center height in the Ferris wheel problem. Q3 is a geometric preview of
  EK 4.4.A.1; connect to the table in the notes.
\end{teachernote}

\end{notesbox}

\end{document}
